\section{Términos clave: índice de palabras clave de este episodio}

\begin{longtable}{p{0.24\textwidth}p{0.34\textwidth}p{0.33\textwidth}}
\toprule
Término & Explicación breve & Papel en este episodio \\
\midrule
Physical AI & IA en el mundo físico, capaz de percibir, decidir y ejecutar acciones & Marco central con el que Xiaopeng (小鹏) pasa de la conducción autónoma a una estrategia de IA. \\
Robotaxi & Servicio de taxis de conducción autónoma & Objetivo de etapa que Liu Xianming (刘先明) quiere que funcione bien en Guangzhou (广州). \\
Software 1.0 & Pila de software dominada por reglas, optimización y lógica escrita a mano & La ruta inicial de la conducción autónoma. \\
Software 2.0 & Pila de software basada en redes neuronales y guiada por datos & Concepto base de la conducción autónoma de extremo a extremo. \\
VLM & Vision-Language Model, modelo de visión y lenguaje & Entiende visión y lenguaje, pero no necesariamente emite una action de forma directa. \\
VLA & Vision-Language-Action, modelo de visión, lenguaje y acción & Dirección importante de la pila tecnológica de Physical AI de Xiaopeng. \\
Language Bottleneck & Cuello de botella intermedio que comprime la información continua de los sensores en token de lenguaje & Capa intermedia clave que Liu Xianming propone eliminar. \\
Self-supervised Learning & Aprendizaje autosupervisado: construye la señal de entrenamiento a partir de los propios datos & Permite aprovechar datos a gran escala y reduce la dependencia del etiquetado manual. \\
Scaling & Mejorar la capacidad con modelos más grandes, más datos y más cómputo & Liu Xianming lo considera la bet clave de Physical AI. \\
Fábrica de modelos en la nube & Entrenar modelos grandes en la nube y luego comprimirlos para desplegarlos en el vehículo & Estructura de ingeniería que conecta el scaling con la producción en serie en el vehículo. \\
Destilación & Usar un modelo grande para entrenar uno pequeño, de modo que el pequeño herede sus capacidades & Medio importante para el despliegue de la nube al vehículo. \\
Cuantización/poda & Reducir la precisión del modelo o recortar parámetros para adaptarlo al hardware & Necesaria para el despliegue en el vehículo y el control de costos. \\
Infra & Infraestructura de datos, entrenamiento, análisis y despliegue & Base que destacan tanto la ruta de Cruise como la de Xiaopeng. \\
Corner Case & Escenarios límite de cola larga & Objeto que el ciclo cerrado de datos de la conducción autónoma debe seguir buscando. \\
\bottomrule
\end{longtable}

\subsection{Resumen de la sección}

Los términos de este episodio giran en torno a una pregunta: cómo llevar la conducción autónoma de la ingeniería basada en reglas a un sistema de Physical AI capaz de hacer scaling. Las palabras clave no son sustantivos aislados, sino una cadena que une entrenamiento, compresión, despliegue y retroalimentación de la producción en serie.
